\documentclass[11pt]{article}
\usepackage[a4paper,margin=1in]{geometry}
\usepackage{amsmath,amssymb,amsthm}
\usepackage[T1]{fontenc}
\usepackage{lmodern}
\usepackage[hidelinks]{hyperref}

\DeclareMathOperator{\SpecFO}{Spec}
\newtheorem{theorem}{Theorem}
\title{The First-Order Spectrum $\{6\}$\\
  \large Disproof of Conjecture 00000009168}
\author{Feiyu\thanks{Independent researcher.}}
\date{26 November 2026}

\begin{document}
\maketitle

A Fagin spectrum is the set of sizes of the nonempty finite models of a
first-order sentence~\cite[\S1]{fagin}. Conjecture 00000009168 characterizes
these spectra as prime-power sets. This characterization requires each
member of each spectrum to be a prime power.

\begin{theorem}
There is a first-order sentence $\varphi$ in the empty signature with
$\SpecFO(\varphi)=\{6\}$. Consequently the prime-power characterization fails.
\end{theorem}

\begin{proof}
For $n\ge 1$, let
\[
  \theta_n =
  \exists x_1\cdots\exists x_n
    \bigwedge_{1\le i<j\le n} (x_i\ne x_j),
  \qquad
  \varphi=\theta_6\land\neg\theta_7 .
\]
A structure satisfies $\theta_n$ exactly when its domain has at least $n$
elements. Thus a model of $\varphi$ has at least six elements and fewer than
seven. Conversely, on the domain $\{0,1,2,3,4,5\}$, its six distinct elements
witness $\theta_6$, and no seven distinct witnesses exist. Equality is
logical, so the domain alone specifies the empty-signature model. Hence
$\SpecFO(\varphi)=\{6\}$.

If $6=p^k$ for a prime $p$ and an integer $k\ge 1$, every prime divisor of
$6$ equals $p$. Both $2$ and $3$ divide $6$, giving a contradiction.
\end{proof}

\paragraph{Lean formalization.}
In \texttt{lean/Submission/Basic.lean}, \texttt{faginSpectrum} quantifies
over finite, nonempty carriers with an interpretation of the language and
uses Mathlib's sentence satisfaction. The sentence \texttt{exactlySix}
conjoins \texttt{Sentence.cardGe} at six with its negation at seven.
\texttt{realize\_exactCard\_iff} proves the cardinality equivalence for every
finite structure, and \texttt{spectrum\_exactlySix} exhibits \texttt{Fin 6}
as a model. \texttt{conjecture\_false} refutes the universal inclusion of
spectra in the set of prime powers. The project uses Lean and Mathlib
\texttt{v4.33.1}.

\begin{thebibliography}{1}
\bibitem{fagin}
R.~Fagin, \emph{Generalized First-Order Spectra and Polynomial-Time
Recognizable Sets}, in Complexity of Computation, SIAM--AMS Proceedings~7
(1974), 43--73.
\href{https://www.cs.umd.edu/users/gasarch/COURSES/752/S26/slides/genspec.pdf}{PDF}.
\end{thebibliography}

\end{document}
